%%%%%%%%%%%%%%%%%%%%%%%%%%%%%%%%%%%%
% Appendix A, Solution and Estimation Details
%%%%%%%%%%%%%%%%%%%%%%%%%%%%%%%%%%%%

% Set equation, figure, table indexing
\renewcommand{\thefigure}{A.\arabic{figure}}
\setcounter{figure}{0}
\renewcommand{\thetable}{A.\arabic{table}}
\setcounter{table}{0}
\renewcommand{\theequation}{A.\arabic{equation}}
\setcounter{equation}{0}
\renewcommand{\thefootnote}{A.\arabic{footnote}}
\setcounter{footnote}{0}
\section*{Data} % (fold)
\label{sec:data-ap}

% =========================
% Implicit association test examples
% =========================

\begin{figure}[H]
\centering
\caption{Examples of an Implicit Association Test}
\label{fig:iatexamples}
\begin{subfigure}{.48\textwidth}
\centering
\includegraphics[width=.9\linewidth]{iatexample1.png}
\end{subfigure}
\centering
%Second graph
\begin{subfigure}{.48\textwidth}
\centering
\includegraphics[width=.9\linewidth]{iatexample2.png}
\end{subfigure}
%Third
\begin{subfigure}{.48\textwidth}
\centering
\includegraphics[width=.9\linewidth]{iatexample3.png}
\end{subfigure}
% Fourth
\begin{subfigure}{.48\textwidth}
\centering
\includegraphics[width=.9\linewidth]{iatexample4.png}
\end{subfigure}
%Fifth
\begin{subfigure}{.48\textwidth}
\centering
\includegraphics[width=.9\linewidth]{iatexample5.png}
\end{subfigure}
\caption*{\footnotesize{Here are a few examples of what a respondent would see on an implicit association test.}}
\end{figure}

\newpage
\pagebreak
% =========================
% Anti-Asian Bias Maps
% =========================


\begin{center}
\begin{figure}[H]
\caption{Maps of State-level Association Test Bias Over Time Measure with Census Division Regional Boundaries}
\label{fig:skiniat-maps}
% first
\begin{subfigure}{.45\textwidth}
\caption{State-level Bias in 2004}
\centering
\includegraphics[width=0.9\linewidth]{2004skinmap.png} 
\label{fig:skiniat-map-2004}
\end{subfigure}
\hfill%
% Second
\begin{subfigure}{.45\textwidth}
\caption{State-level Bias in 2008}
\centering
\includegraphics[width=0.9\linewidth]{2008skinmap.png} 
\label{fig:skiniat-map-2006}
\end{subfigure}
\hfill%
% third
\begin{subfigure}{.45\textwidth}
\caption{State-level Bias in 2012}
\centering
\includegraphics[width=0.9\linewidth]{2012skinmap.png} 
\label{fig:skiniat-map-2008}
\end{subfigure}
\hfill%
% fourth
\begin{subfigure}{.45\textwidth}
\caption{State-level Bias in 2016}
\centering
\includegraphics[width=0.9\linewidth]{2016skinmap.png} 
\label{fig:skiniat-map-2010}
\end{subfigure}

\caption*{\footnotesize{This figure shows the state-level bias index in different years in the sample. The bias units are in standard deviations and ranges from low to high bias. Bias index is constructed following \textcite{lubotskyInterpretationRegressionsMultiple2006}. The data is from the 2004--2021 Harvard's Project Implicit Association Test scores, American National Election Studies (ANES), and state-level hate crimes against Asians. Each panel presents state-level bias during a certain year. The boundaries in black represent the different Census divisions in the United States. Notice how there is a variation across states within a region.}}
\end{figure}
\end{center}

\newpage
\pagebreak

\begin{center}
\begin{figure}[H]
\caption{Maps of State-level Bias 2004--2021 Measure with Census Division Regional Boundaries}
\includegraphics[width=\textwidth]{Average_Skinmap.png} 
\label{fig:iat-map-all}
\caption*{\footnotesize{This figure shows the state-level bias index in the sample from 2004 to 2021. The bias units are in standard deviations and ranges from low to high bias. Bias index is constructed following \textcite{lubotskyInterpretationRegressionsMultiple2006}. The data is from the 2004--2021 Harvard's Project Implicit Association Test scores, American National Election Studies (ANES), and state-level hate crimes against Asians. The boundaries in black represent the different Census divisions in the United States. Notice how there is a variation across states within a region.}}
\end{figure}
\end{center}

\newpage
\pagebreak
% =========================
% Asian identity histograms
% for all generations
% =========================

\begin{center}
\begin{figure}[H]
\caption{Asian Racial Identity}
\includegraphics[width=\textwidth]{histogram_asian_american_race_aggregated.png} 
\label{fig:histogram-all}
\caption*{\footnotesize{This figure shows the distribution of Asian racial identity among respondents. 
The data is aggregated from the 2004--2021 Current Population Survey (CPS). 
The sample includes first-, second-, and third-generation objectively Asian Americans.
Following \textcite{antmanEthnicAttritionObserved2016,antmanEthnicAttritionAssimilation2020}, 
I utilize birthplace information for individuals, parents, and grandparents to create objective Asian ancestry indicators.
}}
\end{figure}
\end{center}

\newpage
\pagebreak

\begin{center}
\begin{figure}[H]
\caption{Asian Racial Identity: First Generation}
\includegraphics[width=\textwidth]{histogram_asian_american_race_firstgen.png} 
\label{fig:histogram-firstgen}
\caption*{\footnotesize{This figure shows the distribution of Asian racial identity among respondents. 
The data is aggregated from the 2004--2021 Current Population Survey (CPS). 
The sample includes first-generation objectively Asian Americans.
Following \textcite{antmanEthnicAttritionObserved2016,antmanEthnicAttritionAssimilation2020}, 
I utilize birthplace information for individuals, parents, and grandparents to create objective Asian ancestry indicators.
A first-generation Asian American is defined as an individual born in an Asian country and is not a US citizen born to US citizen parents abroad.
}}
\end{figure}
\end{center}


\pagebreak
\newpage

\begin{landscape}
\begin{figure}[!htb]
\centering
\caption{Asian Racial Identity: Second Generation}
\label{fig:histogram-secondgen}

% Top row
\begin{subfigure}{.48\textwidth}
\caption{All Second Generation}
\centering
\includegraphics[width=1\linewidth]{histogram_asian_american_race_secondgen.png}
\end{subfigure}
\hfill
\begin{subfigure}{.48\textwidth}
\caption{Second Generation: Asian Fathers-Asian Mothers}
\centering
\includegraphics[width=1\linewidth]{histogram_asian_american_race_secondgen_AA.png}
\end{subfigure}

\vspace{1cm}

% Bottom row
\begin{subfigure}{.48\textwidth}
\caption{Second Generation: Asian Fathers-White Mothers}
\centering
\includegraphics[width=1\linewidth]{histogram_asian_american_race_secondgen_AW.png}
\end{subfigure}
\hfill
\begin{subfigure}{.48\textwidth}
\caption{Second Generation: White Fathers-Asian Mothers}
\centering
\includegraphics[width=1\linewidth]{histogram_asian_american_race_secondgen_WA.png}
\end{subfigure}

\caption*{\footnotesize{This figure shows the distribution of Asian racial identity among respondents. 
The data is aggregated from the 2004--2021 Current Population Survey (CPS). 
The sample includes second-generation objectively Asian Americans.
Following \textcite{antmanEthnicAttritionObserved2016,antmanEthnicAttritionAssimilation2020}, 
I utilize birthplace information for individuals, parents, and grandparents to create objective Asian ancestry indicators.
A second-generation Asian American is defined as a native-born individual with at least one parent born in an Asian country. 
The first panel is for second-generation Asian Americans. The second panel is for second-generation Asian Americans with both parents born in an Asian country. 
The third panel is for second-generation Asian Americans with an Asian father and a White mother. 
The fourth panel is for second-generation Asian Americans with a White father and an Asian mother.
}}
\end{figure}
\end{landscape}

\pagebreak
\newpage

\pagebreak
\newpage

\begin{landscape}
\begin{figure}[!htb]
\centering
\caption{Asian Racial Identity: Third Generation}
\label{fig:histogram-thirdgen}

% Top row - 3 graphs
\begin{subfigure}{.32\textwidth}
\caption{All Third Generation}
\centering
\includegraphics[width=1\linewidth]{histogram_asian_american_race_thirdgen.png}
\end{subfigure}
\hfill
\begin{subfigure}{.32\textwidth}
\caption{Third Generation: One Asian Grandparent}
\centering
\includegraphics[width=1\linewidth]{histogram_asian_american_race_thirdgen_oneasiangran.png}
\end{subfigure}
\hfill
\begin{subfigure}{.32\textwidth}
\caption{Third Generation: Two Asian Grandparents}
\centering
\includegraphics[width=1\linewidth]{histogram_asian_american_race_thirdgen_twoasiangran.png}
\end{subfigure}

\vspace{1cm}

% Bottom row - 2 graphs centered with more space
\hspace*{\fill}
\begin{subfigure}{.4\textwidth}
\caption{Third Generation: Three Asian Grandparents}
\centering
\includegraphics[width=1\linewidth]{histogram_asian_american_race_thirdgen_threeasiangran.png}
\end{subfigure}
\hfill
\begin{subfigure}{.4\textwidth}
\caption{Third Generation: Four Asian Grandparents}
\centering
\includegraphics[width=1\linewidth]{histogram_asian_american_race_thirdgen_fourasiangran.png}
\end{subfigure}
\hspace*{\fill}

\caption*{\footnotesize{This figure shows the distribution of Asian racial identity among respondents. 
The data is aggregated from the 2004--2021 Current Population Survey (CPS). 
The sample includes second-generation objectively Asian Americans.
Following \textcite{antmanEthnicAttritionObserved2016,antmanEthnicAttritionAssimilation2020}, 
I utilize birthplace information for individuals, parents, and grandparents to create objective Asian ancestry indicators.
A third-generation Asian American is defined as a native-born individual with native-born parents and at least one grandparent born in an Asian country. 
The first panel is for third-generation Asian Americans. The second panel is for third-generation Asian Americans with one grandparent born in an Asian country. 
The third panel is for third-generation Asian Americans with two grandparents born in an Asian country. 
The fourth panel is for third-generation Asian Americans with three grandparents born in an Asian country.
The fifth panel is for third-generation Asian Americans with four grandparents born in an Asian country.
}}
\end{figure}
\end{landscape}

\pagebreak
\newpage

\section*{Tables} % (fold)
\label{sec:tabs-ap}

\begin{table}[!h]
\centering\centering
\caption{Subjective Asian Identity and Asian Bias \label{regtab-all-fe}}
\centering
\resizebox{\ifdim\width>\linewidth\linewidth\else\width\fi}{!}{
\begin{threeparttable}
\begin{tabular}[t]{lcccccccc}
\toprule
  & \specialcell{(1) \\ $A_i$} & \specialcell{(2) \\ $A_i$} & \specialcell{(3) \\ $A_i$} & \specialcell{(4) \\ $A_i$} & \specialcell{(5) \\ $A_i$} & \specialcell{(6) \\ $A_i$} & \specialcell{(7) \\ $A_i$} & \specialcell{(8) \\ $A_i$}\\
\midrule
Bias & $-0.04$*** & $-0.14$*** & $-0.02$*** & $-0.02$ & $-0.03$*** & $-0.07$** & $-0.10$*** & $-0.04$\\
 & ($0.01$) & ($0.04$) & ($0.01$) & ($0.03$) & ($0.01$) & ($0.03$) & ($0.03$) & ($0.03$)\\
Female & $-0.01$** & $-0.01$** & $-0.01$** & $-0.01$* & $-0.01$** & $-0.01$* & $-0.01$* & $-0.01$*\\
 & ($0.00$) & ($0.00$) & ($0.00$) & ($0.00$) & ($0.00$) & ($0.00$) & ($0.00$) & ($0.00$)\\
College Graduate: Mother & $0.01$ & $0.01$ & $0.01$ & $0.01$ & $0.01$ & $0.01$ & $0.01$ & $0.01$\\
 & ($0.01$) & ($0.01$) & ($0.01$) & ($0.01$) & ($0.01$) & ($0.01$) & ($0.01$) & \vphantom{2} ($0.01$)\\
College Graduate: Father & $0.01$** & $0.01$* & $0.01$** & $0.01$** & $0.01$** & $0.01$** & $0.01$** & $0.01$**\\
 & ($0.01$) & ($0.01$) & ($0.01$) & ($0.01$) & ($0.01$) & ($0.01$) & ($0.01$) & \vphantom{1} ($0.01$)\\
Both parents Asian & $0.63$*** & $0.63$*** & $0.62$*** & $0.62$*** & $0.63$*** & $0.63$*** & $0.63$*** & $0.62$***\\
 & ($0.03$) & ($0.03$) & ($0.03$) & ($0.03$) & ($0.03$) & ($0.03$) & ($0.03$) & ($0.03$)\\
First Gen & $0.01$ & $0.02$ & $0.02$* & $0.02$** & $0.02$* & $0.02$** & $0.02$* & $0.02$**\\
 & ($0.01$) & ($0.01$) & ($0.01$) & ($0.01$) & ($0.01$) & ($0.01$) & ($0.01$) & ($0.01$)\\
Second Gen & $0.02$ & $0.02$ & $0.01$ & $0.02$ & $0.02$ & $0.02$ & $0.02$ & $0.02$\\
 & ($0.01$) & ($0.01$) & ($0.01$) & ($0.02$) & ($0.02$) & ($0.02$) & ($0.02$) & ($0.02$)\\
\midrule
N & $129078$ & $129078$ & $129078$ & $129078$ & $129078$ & $129078$ & $129078$ & $129078$\\
Region FE &  &  &  &  & X & X &  & \\
Year FE &  & X &  & X &  & X &  & \\
State FE &  &  & X & X &  &  &  & X\\
Year-Region FE &  &  &  &  &  &  & X & X\\
\bottomrule
\multicolumn{9}{l}{\rule{0pt}{1em}* p $<$ 0.1, ** p $<$ 0.05, *** p $<$ 0.01}\\
\end{tabular}
\begin{tablenotes}
\small
\item[1] \footnotesize{I include controls for sex, quartic age, and parental education.}
\item[2] \footnotesize{Standard errors are clustered on the state level.}
\end{tablenotes}
\end{threeparttable}}
\end{table}


\newpage
\pagebreak

\begin{table}[H]
\centering\centering
\caption{Logistic Regression Analysis of Bias and Interracial Marriages \label{regtab-logit-03}}
\centering
\begin{threeparttable}
\begin{tabular}[t]{lccc}
\toprule
\multicolumn{2}{c}{ } & \multicolumn{1}{c}{Asian Men} & \multicolumn{1}{c}{Asian Women} \\
\cmidrule(l{3pt}r{3pt}){3-3} \cmidrule(l{3pt}r{3pt}){4-4}
  & \specialcell{(1) \\ Interracial} & \specialcell{(2) \\ Interracial} & \specialcell{(3) \\ Interracial}\\
\midrule
Bias & $1.47$ & $0.83$ & $1.39$\\
 & ($0.43$) & ($0.33$) & ($0.44$)\\
College Graduate: Wife & $1.42$*** & $1.55$*** & $1.75$***\\
 & ($0.08$) & ($0.19$) & ($0.18$)\\
College Graduate: Husband & $0.94$ & $0.97$ & $0.86$*\\
 & ($0.09$) & ($0.08$) & ($0.07$)\\
\midrule
Observations & $69,800$ & $52,032$ & $60,171$\\
Year $\times$ Region FE & X & X & X\\
\bottomrule
\multicolumn{4}{l}{\rule{0pt}{1em}* p $<$ 0.1, ** p $<$ 0.05, *** p $<$ 0.01}\\
\end{tabular}
\begin{tablenotes}
\small
\item[1] \footnotesize{Each column estimates a logistic regression: $\Pr(interracial_{ist} = 1) = \Lambda(\beta_1 AntiAsianBias_{st} + X_{ist}\pi + \gamma_{rt})$, where $interracial_{ist}$ indicates a couple in which one partner was born in an Asian country and the other was not (labelled Asian husband--White wife or White husband--Asian wife, following the paper's convention), $AntiAsianBias_{st}$ is average anti-Asian bias in state $s$ at time $t$, and $X_{ist}$ includes partner-specific controls (sex, age, education, and years since immigrating to the United States). The unit of observation is a second-generation child; columns (2) and (3) restrict to children of Asian-born fathers and Asian-born mothers, respectively. Coefficients are exponentiated and should be interpreted as odds ratios. Standard errors are clustered on the state level.}
\item[2] \footnotesize{Data source is the 2004-2020 Current Population Survey Data.}
\end{tablenotes}
\end{threeparttable}
\end{table}


\newpage
\pagebreak

% =========================
% Between-Model Tests of AME
% =========================

\begin{table}[H]
\centering
\caption{Between-Model Tests of Average Marginal Effects: Second-Generation AW vs.\ WA \label{tab:between-model-second}}
\centering
\fontsize{10}{12}\selectfont
\begin{threeparttable}
\begin{tabular}[t]{lccc}
\toprule
 & \multicolumn{1}{c}{Asian Only} & \multicolumn{1}{c}{White Only} & \multicolumn{1}{c}{Asian \& White} \\
\cmidrule(l{3pt}r{3pt}){2-2} \cmidrule(l{3pt}r{3pt}){3-3} \cmidrule(l{3pt}r{3pt}){4-4}
\midrule
\multicolumn{4}{l}{\textit{Panel A: AW vs.\ WA}} \\
\addlinespace
Anti-Asian Bias & $-0.024$ & $0.106$** & $-0.082$\\
 & ($0.055$) & ($0.044$) & ($0.062$)\\
\addlinespace
Female & $-0.056$* & $-0.044$ & $0.097$***\\
 & ($0.031$) & ($0.029$) & ($0.028$)\\
\addlinespace
College Graduate: Mother & $0.084$** & $-0.106$** & $0.021$\\
 & ($0.036$) & ($0.050$) & ($0.048$)\\
\addlinespace
College Graduate: Father & $-0.061$** & $0.069$** & $-0.002$\\
 & ($0.027$) & ($0.031$) & ($0.037$)\\
\bottomrule
\multicolumn{4}{l}{\rule{0pt}{1em}* p $<$ 0.1, ** p $<$ 0.05, *** p $<$ 0.01}\\
\end{tabular}
\begin{tablenotes}
\small
\item[1] \footnotesize{Each cell reports the difference in average marginal effects (AME) between the Asian father--White mother (AW) and White father--Asian mother (WA) models, estimated from equation (\ref{eq:multinomial_logit}) using multinomial logit. The difference is computed as $\widehat{\text{AME}}_{\text{AW}} - \widehat{\text{AME}}_{\text{WA}}$.}
\item[2] \footnotesize{Standard errors of the difference are computed from the bootstrap distributions of the two models. $z$-statistics test the null hypothesis that the marginal effects are equal across models.}
\item[3] \footnotesize{I include region $\times$ year fixed effects with controls for quartic age and local Asian population share. Each estimate is the median across 1{,}000 bootstrap resamples, with 95\% percentile confidence intervals.}
\end{tablenotes}
\end{threeparttable}
\end{table}


\pagebreak
\newpage

\begingroup
\fontsize{10}{12}\selectfont
\begin{ThreePartTable}
\begin{TableNotes}
\small
\item[1] \footnotesize{Each cell reports the difference in average marginal effects (AME) between the indicated models, estimated from equation (\ref{eq:multinomial_logit}) using multinomial logit. The difference is computed as $\widehat{\text{AME}}_{\text{Model A}} - \widehat{\text{AME}}_{\text{Model B}}$.}
\item[2] \footnotesize{Standard errors of the difference are computed from the bootstrap distributions of the two models. $z$-statistics test the null hypothesis that the marginal effects are equal across models.}
\item[3] \footnotesize{I include region $\times$ year fixed effects with controls for quartic age and local Asian population share. Each estimate is the median across 1{,}000 bootstrap resamples, with 95\% percentile confidence intervals.}
\item[4] \footnotesize{GP = Grandparent(s). Third-generation individuals are native-born with at least one grandparent born in an Asian country.}
\end{TableNotes}
\begin{longtable}{lccc}
\caption{Between-Model Tests of Average Marginal Effects: Third-Generation by Asian Grandparent Count \label{tab:between-model-third}}\\
\toprule
 & \multicolumn{1}{c}{Asian Only} & \multicolumn{1}{c}{White Only} & \multicolumn{1}{c}{Asian \& White} \\
\cmidrule(l{3pt}r{3pt}){2-2} \cmidrule(l{3pt}r{3pt}){3-3} \cmidrule(l{3pt}r{3pt}){4-4}
\midrule
\endfirsthead
\multicolumn{4}{c}{\tablename\ \thetable{} -- \textit{continued from previous page}}\\
\toprule
 & \multicolumn{1}{c}{Asian Only} & \multicolumn{1}{c}{White Only} & \multicolumn{1}{c}{Asian \& White} \\
\cmidrule(l{3pt}r{3pt}){2-2} \cmidrule(l{3pt}r{3pt}){3-3} \cmidrule(l{3pt}r{3pt}){4-4}
\midrule
\endhead
\midrule
\multicolumn{4}{r}{\textit{Continued on next page}}\\
\endfoot
\bottomrule
\multicolumn{4}{l}{\rule{0pt}{1em}* p $<$ 0.1, ** p $<$ 0.05, *** p $<$ 0.01}\\
\insertTableNotes
\endlastfoot
\multicolumn{4}{l}{\textit{Panel A: One Asian GP vs.\ Two Asian GP}} \\
\addlinespace
Anti-Asian Bias & $0.102$* & $-0.059$ & $-0.047$\\
 & ($0.059$) & ($0.077$) & ($0.084$)\\
\addlinespace
Female & $-0.010$ & $-0.073$ & $0.088$*\\
 & ($0.031$) & ($0.048$) & ($0.050$)\\
\addlinespace
College Graduate: Mother & $-0.048$ & $-0.106$** & $0.160$***\\
 & ($0.032$) & ($0.048$) & ($0.052$)\\
\addlinespace
College Graduate: Father & $-0.021$ & $0.060$ & $-0.041$\\
 & ($0.031$) & ($0.058$) & ($0.061$)\\
\midrule
\multicolumn{4}{l}{\textit{Panel B: One Asian GP vs.\ Three Asian GP}} \\
\addlinespace
Anti-Asian Bias & $0.201$ & $-0.203$*** & $-0.028$\\
 & ($0.152$) & ($0.071$) & ($0.149$)\\
\addlinespace
Female & $0.031$ & $-0.118$*** & $0.087$**\\
 & ($0.031$) & ($0.043$) & ($0.040$)\\
\addlinespace
College Graduate: Mother & $-0.318$*** & $0.198$*** & $0.122$**\\
 & ($0.046$) & ($0.047$) & ($0.051$)\\
\addlinespace
College Graduate: Father & $0.293$*** & $-0.345$*** & $0.057$\\
 & ($0.049$) & ($0.065$) & ($0.062$)\\
\midrule
\multicolumn{4}{l}{\textit{Panel C: Two Asian GP vs.\ Three Asian GP}} \\
\addlinespace
Anti-Asian Bias & $0.098$ & $-0.145$*** & $0.019$\\
 & ($0.156$) & ($0.054$) & ($0.142$)\\
\addlinespace
Female & $0.042$ & $-0.045$ & $-0.000$\\
 & ($0.041$) & ($0.047$) & ($0.045$)\\
\addlinespace
College Graduate: Mother & $-0.270$*** & $0.304$*** & $-0.039$\\
 & ($0.054$) & ($0.055$) & ($0.059$)\\
\addlinespace
College Graduate: Father & $0.314$*** & $-0.405$*** & $0.098$\\
 & ($0.056$) & ($0.060$) & ($0.067$)\\
\end{longtable}
\end{ThreePartTable}
\endgroup


\pagebreak
\newpage
% =========================
% Pooled multinomial logit with
% bias x AW interaction
% =========================

\begin{table}[H]
\centering
\small\setstretch{1.0}
\caption{Pooled Multinomial Logit: Bias Interacted with Parental Configuration, Mixed-Race Second-Generation Children \label{tab:pooled-mnl-bias-x-aw}}
\begin{threeparttable}
\begin{tabular}{lcc}
\toprule
 & White only & Asian and White \\
 & (1) & (2) \\
\midrule
\multicolumn{3}{l}{\textit{Panel A: Coefficients (log odds relative to Asian only)}} \\
Bias & $0.738$** & $0.282$ \\
 & ($0.363$) & ($0.331$) \\
Asian father--White mother (AW) & $-0.615$*** & $-0.675$*** \\
 & ($0.161$) & ($0.153$) \\
Bias $\times$ AW & $0.893$*** & $0.051$ \\
 & ($0.316$) & ($0.360$) \\
\midrule
\multicolumn{3}{l}{\textit{Panel B: Average marginal effects of bias (probability scale)}} \\
Average marginal effect of bias, common-effect model & $0.146$** & $-0.024$ \\
 & ($0.057$) & ($0.067$) \\
Average marginal effect of bias, AW families & $0.249$*** & $-0.056$ \\
 & ($0.068$) & ($0.080$) \\
Average marginal effect of bias, WA families & $0.099$* & $-0.010$ \\
 & ($0.058$) & ($0.070$) \\
Difference, AW minus WA & $0.151$*** & $-0.046$ \\
 & ($0.042$) & ($0.066$) \\
\midrule
\multicolumn{3}{l}{Wald test, Bias $\times$ AW $= 0$ in both equations (clustered): $\chi^2(2) = 10.48$, $p = 0.005$} \\
\multicolumn{3}{l}{Likelihood-ratio test (independent observations): $\chi^2(2) = 106.78$, $p < 0.001$} \\
Observations & \multicolumn{2}{c}{$27{,}048$} \\
State clusters & \multicolumn{2}{c}{$51$} \\
\bottomrule
\end{tabular}
\begin{tablenotes}
\small
\item[1] \footnotesize{The table reports a single multinomial logit, equation (\ref{eq:multinomial_logit}), estimated on pooled second-generation children with one parent born in an Asian country and one born outside Asia, with ``Asian only'' as the reference category. All coefficients are constrained to be equal across Asian father--White mother (AW) and White father--Asian mother (WA) families except that the bias index is interacted with an AW indicator. Controls are sex, mother and father college graduation, a quartic in age, the state Asian population share, and region $\times$ year fixed effects; estimation is weighted by CPS person weights scaled to mean one. 118 observations in 2 region $\times$ year cells with no ``Asian only'' responses are excluded because the cell dummies are not identified. Standard errors in parentheses are clustered at the state level using a sandwich covariance built from the weighted multinomial logit scores. The Wald test uses that clustered covariance; the likelihood-ratio test compares the model with and without the two interaction terms and assumes independent observations. Average marginal effects of a one standard deviation increase in bias are evaluated on the full pooled sample with the AW indicator set to one for every observation and then to zero, so the two rows share the same covariate distribution; the difference is the AW--WA gap with its delta-method standard error. This is a test of whether the bias effect differs by parental configuration, not of the bias effect itself. The first row of Panel B comes from the restricted model with a common bias effect. Conventional standard errors that treat children as independent, the assumption the individual-level bootstrap in the main figures also makes, are smaller: for the common-effect average marginal effects they are 0.018 (White only) and 0.021 (Asian and White) against the clustered 0.057 and 0.067, and for the bias coefficients in Panel A they are 0.128 and 0.114 against 0.363 and 0.331. $^{*}p<0.1$, $^{**}p<0.05$, $^{***}p<0.01$. Data source is the 2004--2021 Current Population Survey.}
\end{tablenotes}
\end{threeparttable}
\end{table}


\pagebreak
\newpage


\section*{Figures}\label{appendix:figs}
% Breakdown: Second Generation by Parental Type (AW and WA)
\begin{center}
\begin{figure}[!htb]
\centering
\caption{Marginal Effects of Key Covariates on Racial Identity Choice by Parental Ancestry (Second-Generation Asian Americans)}
\label{fig:marginal-effects-second-parental}

\begin{subfigure}{.48\textwidth}
\caption{Asian Father-White Mother (AW)}\label{subfig:meaw}
\centering
\includegraphics[width=1\linewidth]{boot_me_second_AW.png}
\end{subfigure}
\hfill
\begin{subfigure}{.48\textwidth}
\caption{White Father-Asian Mother (WA)}\label{subfig:mewa}
\centering
\includegraphics[width=1\linewidth]{boot_me_second_WA.png}
\end{subfigure}

\caption*{\footnotesize{This figure displays average marginal effects of key covariates on the probability of each racial identity outcome, estimated from equation (\ref{eq:multinomial_logit}) separately for second-generation children with an Asian-born father and a mother born outside Asia (panel a) and with a father born outside Asia and an Asian-born mother (panel b). For anti-Asian bias, the marginal effect is the derivative of each predicted probability with respect to the bias index, evaluated for every child at that child's own covariate values and averaged over the estimation sample; for the binary covariates, it is the average change in predicted probability from switching the indicator from 0 to 1 for every child. Models include region $\times$ year fixed effects and controls for a quartic in age and the state Asian population share, and are weighted by CPS person weights. Each point is the median of the average marginal effect across 1{,}000 bootstrap resamples of children, drawn with replacement within region $\times$ year cells; the bars span the 2.5th and 97.5th percentiles of the same resamples and do not account for dependence within state-year cells. Because the two panels come from separately estimated models, differences between them reflect both differences in coefficients and differences in the two groups' covariate distributions; Appendix Table~\ref{tab:pooled-mnl-bias-x-aw} reports a pooled model that holds the covariate distribution fixed.}}
\end{figure}
\end{center}

\pagebreak
\newpage
% Pooled: Second Generation Adults
\begin{center}
\begin{figure}[!htb]
\centering
\caption{Marginal Effects of Key Covariates on Racial Identity Choice (Pooled Mixed-Race Adult Second-Generation Asian Americans)}
\label{fig:marginal-effects-second-adults-pooled}
\includegraphics[width=0.65\linewidth]{logit_boot_me_second_AW_WA_adults_multinom.png}
\caption*{\footnotesize{This figure displays average marginal effects of key covariates on the probability of each racial identity outcome, estimated from equation (\ref{eq:multinomial_logit}) for pooled second-generation adults (ages 18 and older) with one parent born in an Asian country and one parent born outside Asia. For anti-Asian bias, years of education, and log family income, the marginal effect is the derivative of each predicted probability with respect to the covariate, evaluated for every adult at that adult's own covariate values and averaged over the estimation sample; for female, it is the average change in predicted probability from switching the indicator from 0 to 1 for every adult. Models include region $\times$ year fixed effects and controls for a quartic in age and the state Asian population share, and are weighted by CPS person weights. Each point is the median of the average marginal effect across 1{,}000 bootstrap resamples of adults, drawn with replacement within region $\times$ year cells; the bars span the 2.5th and 97.5th percentiles of the same resamples and do not account for dependence within state-year cells.}}
\end{figure}
\end{center}

\pagebreak
\newpage

% % Breakdown: Second Generation Adults by Parental Type (AW and WA)
% \begin{center}
% \begin{figure}[!htb]
% \centering
% \caption{Marginal Effects of Key Covariates on Racial Identity Choice by Parental Ancestry (Adult Second-Generation Asian Americans)}
% \label{fig:marginal-effects-second-adults}

% \begin{subfigure}{.48\textwidth}
% \caption{Asian Father-White Mother (AW)}
% \centering
% \includegraphics[width=1\linewidth]{logit_boot_me_second_AW_adults_multinom.png}
% \end{subfigure}
% \hfill
% \begin{subfigure}{.48\textwidth}
% \caption{White Father-Asian Mother (WA)}
% \centering
% \includegraphics[width=1\linewidth]{logit_boot_me_second_WA_adults_multinom.png}
% \end{subfigure}

% \caption*{\footnotesize{This figure displays the marginal effects of key covariates on the probability of each racial identification outcome, estimated from equation (\ref{eq:multinomial_logit}) using multinomial logit models for second-generation adults with Asian father/White mother (AW) and White father/Asian mother (WA). Marginal effects represent the change in predicted probability for a one-unit increase in the covariate (or discrete change from 0 to 1 for binary variables), holding other variables at their sample means. I include region $\times$ year fixed effects with controls for quartic age and local Asian population share. Each point corresponds to the median marginal effect across 1{,}000 bootstrap resamples, with 95\% percentile confidence intervals. Second-generation Asian Americans are native-born individuals with at least one Asian-born parent.}}
% \end{figure}
% \end{center}

% \pagebreak
% \newpage

% Breakdown: Third Generation by Number of Asian Grandparents
\begin{center}
\begin{figure}[!htb]
\centering
\caption{Marginal Effects of Key Covariates on Racial Identity Choice by Number of Asian Grandparents (Third-Generation Asian Americans)}
\label{fig:marginal-effects-third-grandparental}

\begin{subfigure}{.45\textwidth}
\caption{One Asian Grandparent}\label{subfig:oneasiangrand}
\centering
\includegraphics[width=1\linewidth]{boot_me_third_oneAsian.png}
\end{subfigure}
\hfill
\begin{subfigure}{.45\textwidth}
\caption{Two Asian Grandparents}\label{subfig:twoasiangrand}
\centering
\includegraphics[width=1\linewidth]{boot_me_third_twoAsian.png}
\end{subfigure}

\vspace{0.5cm}

\begin{subfigure}{.45\textwidth}
\caption{Three Asian Grandparents}
\centering
\includegraphics[width=1\linewidth]{boot_me_third_threeAsian.png}
\end{subfigure}

\caption*{\footnotesize{This figure displays average marginal effects of key covariates on the probability of each racial identity outcome, estimated from equation (\ref{eq:multinomial_logit}) separately for third-generation children with one, two, and three Asian-born grandparents. For anti-Asian bias, the marginal effect is the derivative of each predicted probability with respect to the bias index, evaluated for every child at that child's own covariate values and averaged over the estimation sample; for the binary covariates, it is the average change in predicted probability from switching the indicator from 0 to 1 for every child. Models include region $\times$ year fixed effects and controls for a quartic in age and the state Asian population share, and are weighted by CPS person weights. Each point is the median of the average marginal effect across 1{,}000 bootstrap resamples of children, drawn with replacement within region $\times$ year cells; the bars span the 2.5th and 97.5th percentiles of the same resamples.}}
\end{figure}
\end{center}

\pagebreak
\newpage
% =========================
% Between-Model Tests: Second Generation (AW vs. WA)
% =========================
\begin{center}
\begin{figure}[!htb]
\centering
\caption{Between-Model Differences in Average Marginal Effects: Second-Generation AW vs.\ WA}
\label{fig:between-model-second}

\includegraphics[width=0.85\linewidth]{between_model_second_AW_WA.png}

\caption*{\footnotesize{This figure plots the difference in average marginal effects (AME) between the Asian father--White mother (AW) and White father--Asian mother (WA) multinomial logit models, estimated from equation (\ref{eq:multinomial_logit}). Each point represents $\widehat{\text{AME}}_{\text{AW}} - \widehat{\text{AME}}_{\text{WA}}$ for a given covariate and racial identity outcome. Filled markers indicate statistical significance at the 5\% level. Error bars show 95\% confidence intervals constructed from the bootstrap standard errors of the difference. I include region $\times$ year fixed effects with controls for quartic age and local Asian population share. Each estimate is the median across 1{,}000 bootstrap resamples.}}
\end{figure}
\end{center}

\pagebreak
\newpage

% =========================
% Between-Model Tests: Third Generation by Grandparent Count
% =========================
\begin{center}
\begin{figure}[!htb]
\centering
\caption{Between-Model Differences in Average Marginal Effects: Third-Generation by Asian Grandparent Count}
\label{fig:between-model-third}

\begin{subfigure}{.48\textwidth}
\caption{One vs.\ Two Asian Grandparents}
\centering
\includegraphics[width=1\linewidth]{between_model_third_1v2GP.png}
\end{subfigure}
\hfill
\begin{subfigure}{.48\textwidth}
\caption{One vs.\ Three Asian Grandparents}
\centering
\includegraphics[width=1\linewidth]{between_model_third_1v3GP.png}
\end{subfigure}

\vspace{0.5cm}

\begin{subfigure}{.48\textwidth}
\caption{Two vs.\ Three Asian Grandparents}
\centering
\includegraphics[width=1\linewidth]{between_model_third_2v3GP.png}
\end{subfigure}

\caption*{\footnotesize{This figure plots the difference in average marginal effects (AME) between third-generation multinomial logit models estimated on samples defined by the number of Asian grandparents. Each point represents the difference in AME between the indicated models for a given covariate and racial identity outcome. Filled markers indicate statistical significance at the 5\% level. Error bars show 95\% confidence intervals constructed from the bootstrap standard errors of the difference. I include region $\times$ year fixed effects with controls for quartic age and local Asian population share. Each estimate is the median across 1{,}000 bootstrap resamples.}}
\end{figure}
\end{center}

% \pagebreak
% \newpage

% \begin{center}
% \begin{figure}[!htb]
% \centering
% \caption{Multinomial Logit Model: Predicted Probabilities of Racial Identity Choice by Key Covariates (Second-Generation Asian Americans with Asian Fathers and White Mothers)}
% \label{fig:pp-second-aw}

% \begin{subfigure}{.48\textwidth}
% \caption{Anti-Asian Bias}
% \centering
% \includegraphics[width=1\linewidth]{boot_pp_second_AW_value.png}
% \end{subfigure}
% \hfill
% \begin{subfigure}{.48\textwidth}
% \caption{Female}
% \centering
% \includegraphics[width=1\linewidth]{boot_pp_second_AW_Female.png}
% \end{subfigure}

% \vspace{0.5cm}

% \begin{subfigure}{.48\textwidth}
% \caption{College Graduate: Mother}
% \centering
% \includegraphics[width=1\linewidth]{boot_pp_second_AW_MomGradCollege.png}
% \end{subfigure}
% \hfill
% \begin{subfigure}{.48\textwidth}
% \caption{College Graduate: Father}
% \centering
% \includegraphics[width=1\linewidth]{boot_pp_second_AW_DadGradCollege.png}
% \end{subfigure}

% \caption*{\footnotesize{This figure shows predicted probabilities from estimating equation (\ref{eq:multinomial_logit}) using a multinomial logit model for the subsample of second-generation Asian Americans with Asian fathers and White mothers. The model estimates the probability of choosing ``Asian only'', ``White only'', or ``Asian and White'' racial identification as a function of anti-Asian bias, gender, and parental education. I include region $\times$ year fixed effects with controls for quartic age and local Asian population share. All other variables are held at their sample means. Reported curves depict the median predicted probability across 1{,}000 bootstrap resamples, and 95\% confidence intervals follow the corresponding bootstrap percentile limits.}}
% \end{figure}
% \end{center}

% \begin{center}
% \begin{figure}[!htb]
% \centering
% \caption{Multinomial Logit Model: Predicted Probabilities of Racial Identity Choice by Key Covariates (Second-Generation Asian Americans with White Fathers and Asian Mothers)}
% \label{fig:pp-second-wa}

% \begin{subfigure}{.48\textwidth}
% \caption{Anti-Asian Bias}
% \centering
% \includegraphics[width=1\linewidth]{boot_pp_second_WA_value.png}
% \end{subfigure}
% \hfill
% \begin{subfigure}{.48\textwidth}
% \caption{Female}
% \centering
% \includegraphics[width=1\linewidth]{boot_pp_second_WA_Female.png}
% \end{subfigure}

% \vspace{0.5cm}

% \begin{subfigure}{.48\textwidth}
% \caption{College Graduate: Mother}
% \centering
% \includegraphics[width=1\linewidth]{boot_pp_second_WA_MomGradCollege.png}
% \end{subfigure}
% \hfill
% \begin{subfigure}{.48\textwidth}
% \caption{College Graduate: Father}
% \centering
% \includegraphics[width=1\linewidth]{boot_pp_second_WA_DadGradCollege.png}
% \end{subfigure}

% \caption*{\footnotesize{This figure shows predicted probabilities from estimating equation (\ref{eq:multinomial_logit}) using a multinomial logit model for the subsample of second-generation Asian Americans with White fathers and Asian mothers. The model estimates the probability of choosing ``Asian only'', ``White only'', or ``Asian and White'' racial identification as a function of anti-Asian bias, gender, and parental education. I include region $\times$ year fixed effects with controls for quartic age and local Asian population share. All other variables are held at their sample means. Reported curves depict the median predicted probability across 1{,}000 bootstrap resamples, and 95\% confidence intervals follow the corresponding bootstrap percentile limits.}}
% \end{figure}
% \end{center}


% % =========================
% % 41: Multinomial Model - Predicted Probabilities (Adults, AW/WA)
% % =========================
% \begin{center}
% \begin{figure}[!htb]
% \centering
% \caption{Multinomial Logit Model: Predicted Probabilities of Racial Identity Choice by Anti-Asian Bias and Gender (Second Generation Adults, AW/WA)}
% \label{fig:pp-secondgen-adults-aw-wa}

% % Top row: Anti-Asian Bias and Female
% \begin{subfigure}{.48\textwidth}
% \caption{Anti-Asian Bias (AW)}
% \centering
% \includegraphics[width=1\linewidth]{logit_boot_pp_second_AW_adults_multinom_value.png}\label{subfig:pp-secgen-aw-bias}
% \end{subfigure}
% \hfill
% \begin{subfigure}{.48\textwidth}
% \caption{Anti-Asian Bias (WA)}
% \centering
% \includegraphics[width=1\linewidth]{logit_boot_pp_second_WA_adults_multinom_value.png}\label{subfig:pp-secgen-wa-bias}
% \end{subfigure}

% \vspace{0.5cm}

% % Second row: Female
% \begin{subfigure}{.48\textwidth}
% \caption{Female (AW)}
% \centering
% \includegraphics[width=1\linewidth]{logit_boot_pp_second_AW_adults_multinom_Female.png}\label{subfig:pp-secgen-aw-female}
% \end{subfigure}
% \hfill
% \begin{subfigure}{.48\textwidth}
% \caption{Female (WA)}
% \centering
% \includegraphics[width=1\linewidth]{logit_boot_pp_second_WA_adults_multinom_Female.png}\label{subfig:pp-secgen-wa-female}
% \end{subfigure}

% \caption*{\footnotesize{Predicted probabilities from multinomial logit models for second-generation adults with Asian father/White mother (AW) and White father/Asian mother (WA). Each panel shows the probability of choosing each racial identity as a function of anti-Asian bias or gender, holding other covariates at their means. Curves plot the median predicted probability across 1{,}000 bootstrap resamples, with 95\% confidence intervals given by the bootstrap percentile bounds.}}
% \end{figure}
% \end{center}

% \begin{center}
% \begin{figure}[!htb]
% \centering
% \caption{Multinomial Logit Model: Predicted Probabilities of Racial Identity Choice by Years of Education and Income (Second Generation Adults, AW/WA)}
% \label{fig:pp-secondgen-adults-aw-wa-2}

% % Top row: EducYears
% \begin{subfigure}{.48\textwidth}
% \caption{Years of Education (AW)}
% \centering
% \includegraphics[width=1\linewidth]{logit_boot_pp_second_AW_adults_multinom_EducYears.png}\label{subfig:pp-secgen-aw-EducYears}
% \end{subfigure}
% \hfill
% \begin{subfigure}{.48\textwidth}
% \caption{Years of Education (WA)}
% \centering
% \includegraphics[width=1\linewidth]{logit_boot_pp_second_WA_adults_multinom_EducYears.png}\label{subfig:pp-secgen-wa-EducYears}
% \end{subfigure}

% \vspace{0.5cm}

% % Second row: lnftotval
% \begin{subfigure}{.48\textwidth}
% \caption{Log Family Income (AW)}
% \centering
% \includegraphics[width=1\linewidth]{logit_boot_pp_second_AW_adults_multinom_lnftotval.png}\label{subfig:pp-secgen-aw-lnftotval}
% \end{subfigure}
% \hfill
% \begin{subfigure}{.48\textwidth}
% \caption{Log Family Income (WA)}
% \centering
% \includegraphics[width=1\linewidth]{logit_boot_pp_second_WA_adults_multinom_lnftotval.png}\label{subfig:pp-secgen-wa-lnftotval}
% \end{subfigure}

% \caption*{\footnotesize{Predicted probabilities from multinomial logit models for second-generation adults with Asian father/White mother (AW) and White father/Asian mother (WA). Each panel shows the probability of choosing each racial identity as a function of total years of education or log family income, holding other covariates at their means. Curves plot the median predicted probability across 1{,}000 bootstrap resamples, with 95\% confidence intervals given by the bootstrap percentile bounds.}}
% \end{figure}
% \end{center}

% % =========================
% % Multinomial Model - Predicted Probabilities Third Gens
% % =========================

% \begin{center}
% \begin{figure}[!htb]
% \centering
% \caption{Multinomial Logit Model: Predicted Probabilities of Racial Identity Choice by Key Covariates (Third-Generation Asian Americans with One Asian Grandparent)}
% \label{fig:pp-third-one}

% \begin{subfigure}{.48\textwidth}
% \caption{Anti-Asian Bias}
% \centering
% \includegraphics[width=1\linewidth]{boot_pp_third_oneAsian_value.png}
% \end{subfigure}
% \hfill
% \begin{subfigure}{.48\textwidth}
% \caption{Female}
% \centering
% \includegraphics[width=1\linewidth]{boot_pp_third_oneAsian_Female.png}
% \end{subfigure}

% \vspace{0.5cm}

% \begin{subfigure}{.48\textwidth}
% \caption{College Graduate: Mother}
% \centering
% \includegraphics[width=1\linewidth]{boot_pp_third_oneAsian_MomGradCollege.png}
% \end{subfigure}
% \hfill
% \begin{subfigure}{.48\textwidth}
% \caption{College Graduate: Father}
% \centering
% \includegraphics[width=1\linewidth]{boot_pp_third_oneAsian_DadGradCollege.png}
% \end{subfigure}

% \caption*{\footnotesize{This figure shows predicted probabilities from estimating equation (\ref{eq:multinomial_logit}) using a multinomial logit model for the subsample of third-generation Asian Americans with one Asian grandparent. The model estimates the probability of choosing ``Asian only'', ``White only'', or ``Asian and White'' racial identification as a function of anti-Asian bias, gender, and parental education. I include region $\times$ year fixed effects with controls for quartic age and local Asian population share. All other variables are held at their sample means. Reported curves depict the median predicted probability across 1{,}000 bootstrap resamples, with 95\% percentile confidence intervals from the same draws.}}
% \end{figure}
% \end{center}

% \begin{center}
% \begin{figure}[!htb]
% \centering
% \caption{Multinomial Logit Model: Predicted Probabilities of Racial Identity Choice by Key Covariates (Third-Generation Asian Americans with Two Asian Grandparents)}\label{fig:pp-third-two}

% \begin{subfigure}{.48\textwidth}
% \caption{Anti-Asian Bias}
% \centering
% \includegraphics[width=1\linewidth]{boot_pp_third_twoAsian_value.png}
% \end{subfigure}
% \hfill
% \begin{subfigure}{.48\textwidth}
% \caption{Female}
% \centering
% \includegraphics[width=1\linewidth]{boot_pp_third_twoAsian_Female.png}
% \end{subfigure}

% \vspace{0.5cm}

% \begin{subfigure}{.48\textwidth}
% \caption{College Graduate: Mother}
% \centering
% \includegraphics[width=1\linewidth]{boot_pp_third_twoAsian_MomGradCollege.png}
% \end{subfigure}
% \hfill
% \begin{subfigure}{.48\textwidth}
% \caption{College Graduate: Father}
% \centering
% \includegraphics[width=1\linewidth]{boot_pp_third_twoAsian_DadGradCollege.png}
% \end{subfigure}

% \caption*{\footnotesize{This figure shows predicted probabilities from estimating equation (\ref{eq:multinomial_logit}) using a multinomial logit model for the subsample of third-generation Asian Americans with two Asian grandparents. The model estimates the probability of choosing ``Asian only'', ``White only'', or ``Asian and White'' racial identification as a function of anti-Asian bias, gender, and parental education. I include region $\times$ year fixed effects with controls for quartic age and local Asian population share. All other variables are held at their sample means. Reported curves depict the median predicted probability across 1{,}000 bootstrap resamples, with 95\% percentile confidence intervals from the same draws.}}
% \end{figure}
% \end{center}

\pagebreak
\newpage

\begin{figure}[!htb]
    \centering
    \caption{Scatter Plot of Proportion Subjectively Asian on Bias}
    \label{scatter-plot-1}
    \begin{subfigure}{.9\textwidth}
    \caption{Year < 2015}
    \centering
    \includegraphics[width=.9\linewidth]{scatter-plot-bias-Asian-less2015.png}
    \end{subfigure}
    \centering
    %Second graph
    \begin{subfigure}{.9\textwidth}
    \caption{Year $\geq$ 2015}
    \centering
    \includegraphics[width=.9\linewidth]{scatter-plot-bias-Asian-great2015.png}
    \end{subfigure}
    \caption*{\footnotesize{Here are two scatter plots showing the relationship between bias and subjective Asian population in a state. Each dot represents a state in a certain year. Percent subjectively Asian = $\frac{\# \text{Asian}}{\text{Population}}$ \\
    \emph{Source.} 2004--2021 Current Population Survey.}}
\end{figure}   

\newpage
\pagebreak
\FloatBarrier

\subsection*{Using \textcite{lubotskyInterpretationRegressionsMultiple2006} to Construct Bias Index} % (fold)
\label{sub:lw-bias-appendix}
In \textcite{lubotskyInterpretationRegressionsMultiple2006}, the authors propose a method to reduce measurement error in proxies by constructing a composite index. The Lubotsky-Wittenberg (henceforth LW) consider a model where a covariate is unobserved. Therefore, they use two proxies in its place, which will have measurement error. Thus, the LW method allows researchers to use two proxies that are error-ridden. 

LW consider a setup with the following model:

\begin{align*}
    y &= \alpha + \beta x^* + \epsilon \\
    x_1 &= x^* + \mu_1 \\
    x_2 &= x^* + \mu_2
\end{align*}

Where \(x_{i}^{*}\) is the unobserved covariate, \(x_{1i}\) and \(x_{2i}\) are the proxies, and the measurement errors \(\mu_1\) and \(\mu_2\) are assumed to be classical and allowed to covary. The covariance matrix of the errors is given by:

\[
\Sigma = \begin{bmatrix}
    \sigma_1^2 & \sigma_{12} \\
    \sigma_{12} & \sigma_2^2
\end{bmatrix}
\]

Replacing the unobserved \(x^*\) with \(x_1\) or \(x_2\) yields the following expectations of the OLS estimates:

\[
\mathbb{E} \left[ \hat{\beta}_1 \right] = \beta \frac{\sigma_{x^*}^2}{\sigma_{x^*}^2 + \sigma_1^2} \quad ; \quad \mathbb{E} \left[ \hat{\beta}_2 \right] = \beta \frac{\sigma_{x^*}^2}{\sigma_{x^*}^2 + \sigma_2^2}
\]

Both estimates are biased; the one with the smaller variance of the measurement error being less biased.

LW then propose defining a new proxy \(x_3\) as a weighted average of \(x_1\) and \(x_2\):

\[
x_3 = \lambda x_1 + (1 - \lambda) x_2
\]

To minimize the attenuation bias in the OLS estimate of \(\beta\), they solve for the optimal value of \(\lambda\):

\[
\lambda^* = \frac{\sigma_2^2 - \sigma_{12}}{\sigma_1^2 + \sigma_2^2 - 2\sigma_{12}}
\]

This optimal value of \(\lambda\) is not directly useful because the variances of the measurement errors and their covariance are unobserved. However, if you estimate a bivariate regression using OLS (i.e., regress \(y\) on \(x_1\) and \(x_2\)), then the expectation of the sum of the two coefficient estimates is identical to the expectation of the OLS coefficient estimate on \(x_3\) in a univariate regression using the optimal choice of \(\lambda\):

\[
\mathbb{E} \left[ \hat{\beta}_1 + \hat{\beta}_2 \right] = \mathbb{E} \left[ \hat{\beta}_{x_3} \right]
\]

Thus, OLS produces an estimate of \(\beta\) with the least bias by optimally combining the information in \(x_1\) and \(x_2\).
